% ============================================================================
\chapter{模形式、椭圆曲线与模曲线}
\label{ch:forms-curves}
% ============================================================================
% 对应 Diamond \& Shurman Chapter 1

% --- 章导言：从一个具体函数出发，引起好奇心 ---

在正式定义之前，让我们先认识一个具体的函数。取
\[
  E_4(\tau) = 1 + 240\sum_{n=1}^{\infty} \sigma_3(n)\, q^n,
  \qquad q = e^{2\pi i\tau},
\]
其中 $\sigma_3(n) = \sum_{d \mid n} d^3$ 是因子的立方和。
前几项展开为
\[
  E_4(\tau) = 1 + 240q + 2160q^2 + 6720q^3 + \cdots
\]
这个级数收敛的区域是\textbf{上半平面} $\HH = \{\tau \in \C \mid \im(\tau) > 0\}$
（因为此时 $|q| = e^{-2\pi \im(\tau)} < 1$）。

$E_4$ 有两个令人惊叹的对称性：
\begin{enumerate}
  \item \textbf{平移不变}：$E_4(\tau + 1) = E_4(\tau)$
    ——这很"显然"，因为 $e^{2\pi i(\tau+1)} = e^{2\pi i\tau}$，
    所以 $q$ 没变。
  \item \textbf{反演对称}：$E_4(-1/\tau) = \tau^4 \cdot E_4(\tau)$
    ——这就不那么明显了！把 $\tau$ 替换成 $-1/\tau$，
    整个函数居然只差一个 $\tau^4$ 的因子。
\end{enumerate}
这些性质的证明并不平凡：反演对称性需要对求和指标做巧妙的重排
（Hecke 的技巧），而 $q$-展开的推导则用到 Poisson 求和公式。
我们将在\cref{ch:eisenstein}中给出完整证明。

这两条对称性把 $E_4$ 变成了一个非常特殊的函数，叫做\textbf{模形式}（modular form）。
它的"权"（weight）是 $4$——就是那个 $\tau^4$ 的指数。

本章的任务就是回答三个问题：
\begin{itemize}
  \item \textbf{对称性从哪来？}
    平移 $\tau \mapsto \tau+1$ 和反演 $\tau \mapsto -1/\tau$
    其实是某个群的两个生成元——它就是\textbf{模群} $\SL_2(\Z)$。
  \item \textbf{模形式空间有多大？}
    满足这些对称性的函数居然只构成一个\textbf{有限维}向量空间。
    例如权 $12$ 的模形式空间只有 $2$ 维！
  \item \textbf{模形式跟椭圆曲线有什么关系？}
    上半平面在模群作用下的商空间参数化了所有椭圆曲线。
    这个联系最终通向费马大定理。
\end{itemize}


% ============================================================================

\section{模群与上半平面}
\label{sec:modular-group-action}
% ============================================================================

上面我们观察到 $E_4$ 在平移 $\tau \mapsto \tau + 1$ 和反演 $\tau \mapsto -1/\tau$
下有特殊行为。这两个变换可以组合出更多变换，它们构成一个群。
让我们把这个群精确地描述出来。

\begin{definition}[上半平面]
\label{def:upper-half-plane}
\textbf{上半平面}\index{上半平面 (upper half-plane)}\index{H@$\mathbb{H}$}（upper half-plane）定义为
\[
  \HH = \{ \tau \in \C \mid \im(\tau) > 0 \}.
\]
\end{definition}

\begin{definition}[模群]
\label{def:modular-group}
\index{模群 (modular group)}\index{SL2Z@$\SL_2(\mathbb{Z})$}
\textbf{模群} $\SL_2(\Z)$ 是所有行列式为 $1$ 的 $2 \times 2$ 整数矩阵组成的群：
\[
  \SL_2(\Z) = \left\{
    \gamma = \begin{pmatrix} a & b \\ c & d \end{pmatrix}
    \;\middle|\; a, b, c, d \in \Z,\; ad - bc = 1
  \right\}.
\]
它通过\textbf{莫比乌斯变换}\index{莫比乌斯变换 (Möbius transformation)}（Möbius transformation）作用在 $\HH$ 上：
\[
  \gamma(\tau) = \frac{a\tau + b}{c\tau + d}.
\]
\end{definition}

\begin{proposition}
\label{prop:mobius-preserves-H}
莫比乌斯变换保持上半平面不变：若 $\tau \in \HH$，$\gamma \in \SL_2(\Z)$，
则 $\gamma(\tau) \in \HH$。具体地，
\[
  \im(\gamma(\tau)) = \frac{\im(\tau)}{|c\tau + d|^2}.
\]
\end{proposition}

\begin{proof}
设 $\gamma = \begin{pmatrix} a & b \\ c & d \end{pmatrix}$，$\tau = x + iy$。
\begin{align}
  \gamma(\tau)
  &= \frac{a\tau + b}{c\tau + d}
  = \frac{(a\tau + b)\overline{(c\tau + d)}}{|c\tau + d|^2}.
\end{align}
分子的虚部为
\begin{align}
  \im\bigl((a\tau + b)\overline{(c\tau + d)}\bigr)
  &= \im\bigl((a\tau + b)(\bar c \bar\tau + \bar d)\bigr) \\
  &= \im\bigl(ac|\tau|^2 + ad\tau + bc\bar\tau + bd\bigr) \\
  &= (ad - bc)\im(\tau) = \im(\tau).
\end{align}
最后一步用了 $ad - bc = 1$ 以及 $\im(ad\tau + bc\bar\tau) = (ad - bc)y$。
故 $\im(\gamma(\tau)) = \im(\tau)/|c\tau + d|^2 > 0$。
\end{proof}

\textbf{为什么用 $\SL_2$ 而非 $\GL_2$？}
$\GL_2(\Z)$ 包含行列式为 $\pm 1$ 的矩阵。行列式为 $-1$ 的矩阵
（如 $\begin{pmatrix} 1 & 0 \\ 0 & -1 \end{pmatrix}$）对应变换
$\tau \mapsto -\tau$，把上半平面映到\emph{下}半平面。由上面的证明，
$\im(\gamma(\tau)) = (ad - bc)\im(\tau)/|c\tau + d|^2$，
只有 $ad - bc = +1$（即 $\det \gamma = 1$）时虚部才保持正性。

\textbf{$-I$ 的作用与 $\PSL_2(\Z)$。}
$-I = \begin{pmatrix} -1 & 0 \\ 0 & -1 \end{pmatrix}$ 作用平凡：
$(-I)(\tau) = \frac{-\tau}{-1} = \tau$。
也就是说，$\SL_2(\Z)$ 中不同的矩阵可能对 $\HH$ 做同样的事
（$\gamma$ 和 $-\gamma$ 永远给出相同的变换）。
把这种"冗余"消除掉，就得到
$\PSL_2(\Z) = \SL_2(\Z)/\{\pm I\}$——
在 $\PSL_2(\Z)$ 中，不同的元素一定对应不同的变换。

\begin{theorem}[$\SL_2(\Z)$ 的生成元]
\label{thm:sl2z-generators}
$\SL_2(\Z)$ 由以下两个矩阵生成：
\[
  S = \begin{pmatrix} 0 & -1 \\ 1 & 0 \end{pmatrix}, \quad
  T = \begin{pmatrix} 1 & 1 \\ 0 & 1 \end{pmatrix}.
\]
它们对应的莫比乌斯变换为
\[
  S(\tau) = -\frac{1}{\tau}, \qquad T(\tau) = \tau + 1.
\]
注意 $S$ 在 $\SL_2(\Z)$ 中阶为 $4$（$S^2 = -I$，$S^4 = I$），
在 $\PSL_2(\Z)$ 中阶为 $2$。
\end{theorem}

\textbf{为什么这个定理很重要？}
定义模形式时，我们要求 $f(\gamma(\tau)) = (c\tau+d)^k f(\tau)$ 对
\emph{所有} $\gamma \in \SL_2(\Z)$ 成立——
这是无穷多个条件。
但\cref{thm:sl2z-generators} 告诉我们：$\SL_2(\Z)$ 中的每个矩阵都是 $S$ 和 $T$ 的乘积，
所以只要验证两条：
\[
  f(\tau + 1) = f(\tau) \quad \text{和} \quad f(-1/\tau) = \tau^k f(\tau),
\]
无穷多个条件就自动都满足了。这把"验证对称性"的工作量从无穷压缩到了两条。

\textbf{$S$ 与 $T$ 的几何直觉。}
\label{rem:S-T-geometry}
\begin{itemize}
  \item $T(\tau) = \tau + 1$：水平平移。将竖直带
    $\{-1/2 < \Re(\tau) \leq 1/2\}$ 左右"复制粘贴"铺满整个上半平面。
  \item $S(\tau) = -1/\tau$：关于单位圆的反演（加上实轴的反射）。
    关键性质：$|S(\tau)| = 1/|\tau|$，因此
    \begin{itemize}
      \item 单位圆 $|\tau| = 1$ 被映到\emph{自身}；
      \item 单位圆外部（$|\tau| > 1$）映到内部（$|S(\tau)| < 1$），反之亦然。
    \end{itemize}
    这解释了基本域 $\mathcal{F}$ 的形状：
    底边是单位圆弧（$S$ 的分界线），两侧是 $\Re(\tau) = \pm 1/2$
    （$T$ 的分界线）。
\end{itemize}

\begin{proof}
需要证明任意 $\gamma = \begin{pmatrix} a & b \\ c & d \end{pmatrix} \in \SL_2(\Z)$
可以写成 $S$ 和 $T$ 的乘积。

我们对 $|c|$ 归纳。

\textbf{基础情形 $c = 0$}：由 $ad = 1$，得 $a = d = \pm 1$。
\[
  \begin{pmatrix} 1 & b \\ 0 & 1 \end{pmatrix} = T^b, \quad
  \begin{pmatrix} -1 & -b \\ 0 & -1 \end{pmatrix} = S^2 T^{-b}.
\]
（因为 $S^2 = -I$。）

\textbf{归纳步骤}：设 $c \neq 0$。用带余除法：$a = qc + r$，$0 \leq r < |c|$。
注意
\[
  T^{-q} \gamma
  = \begin{pmatrix} 1 & -q \\ 0 & 1 \end{pmatrix}
    \begin{pmatrix} a & b \\ c & d \end{pmatrix}
  = \begin{pmatrix} a - qc & b - qd \\ c & d \end{pmatrix}
  = \begin{pmatrix} r & b-qd \\ c & d \end{pmatrix}.
\]
再左乘 $S$：
\[
  S \cdot T^{-q} \gamma
  = \begin{pmatrix} 0 & -1 \\ 1 & 0 \end{pmatrix}
    \begin{pmatrix} r & b-qd \\ c & d \end{pmatrix}
  = \begin{pmatrix} -c & -d \\ r & b-qd \end{pmatrix}.
\]
新矩阵的左下角元素为 $r$，且 $0 \leq r < |c|$。
由归纳假设，$ST^{-q}\gamma$ 可用 $S, T$ 表示，故 $\gamma$ 也可以。
\end{proof}


% ============================================================================

\section{基本域}
\label{sec:fundamental-domain}
% ============================================================================

$\SL_2(\Z)$ 包含无穷多个矩阵，每一个都把上半平面映到自身。
自然的问题是：上半平面中哪些点在 $\SL_2(\Z)$ 看来是"本质上相同的"？

\textbf{先想想我们想得到什么。}
如果 $\tau$ 和 $\tau'$ 满足 $\tau' = \gamma(\tau)$ 对某个 $\gamma \in \SL_2(\Z)$，
那么任何模形式 $f$ 在这两个点的值满足
$f(\tau') = (c\tau+d)^k f(\tau)$——它们"本质上一样"。
这意味着 $\HH$ 中有大量"冗余"：无穷多个点实际上代表同一条椭圆曲线。

\textbf{基本域的作用}就是消除这种冗余：
从每个轨道 $\SL_2(\Z) \cdot \tau$ 中恰好挑一个"标准代表元"，
所有代表元合在一起构成的区域就是基本域。
有了基本域，我们就可以把"$\HH$ 上的问题"化简成"基本域这个小区域上的问题"，
从而把无限复杂的空间变成有限大小的对象来分析。

\textbf{实际用途：}
基本域在后面的维数公式（\cref{ch:dimension}）中至关重要。
计算 $\dim \Mk_k(\Gamma)$ 时，需要在商空间 $\Gamma \backslash \HH^*$ 上用 Riemann-Hurwitz 公式，
而基本域就是这个商空间的一个具体几何实现。

这需要\textbf{群作用}\index{群作用 (group action)}（group action）的语言。群 $G$ \textbf{作用在}集合 $X$ 上，
是指一个映射 $G \times X \to X$（记 $(g, x) \mapsto g \cdot x$），满足
$e \cdot x = x$ 和 $(g_1 g_2) \cdot x = g_1 \cdot (g_2 \cdot x)$。
在我们的情况下，$G = \SL_2(\Z)$，$X = \HH$，作用就是莫比乌斯变换。

两个关键概念：
\begin{itemize}
  \item \textbf{轨道}\index{轨道 (orbit)}（orbit）：$G \cdot \tau = \{\gamma(\tau) \mid \gamma \in G\}$，
    即 $\tau$ 能被 $G$ 映到的所有点。同一轨道中的点视为"等价"的。
  \item \textbf{基本域}：从每个轨道中恰好选一个代表元，
    这些代表元组成的区域就是基本域——上半平面的一个"最小切片"。
\end{itemize}

\begin{example}[轨道是什么样的]
\label{ex:orbit-of-2i}
取 $\tau = 2i$。它的轨道 $\SL_2(\Z) \cdot (2i)$ 包含哪些点？
\begin{itemize}
  \item $T^n(2i) = 2i + n$，$n \in \Z$
    ——整条水平线 $\{\ldots, -1+2i, \; 2i, \; 1+2i, \; 2+2i, \ldots\}$。
  \item $S(2i) = -1/(2i) = i/2$
    ——跳到了更低的位置。
  \item $T(S(2i)) = i/2 + 1 = 1 + i/2$，$T^{-1}(S(2i)) = -1 + i/2$，……
  \item $S(T(2i)) = S(1+2i) = \frac{-1}{1+2i} = \frac{-(1-2i)}{5}
    = -\frac{1}{5} + \frac{2}{5}i$
    ——又跳到了别的地方。
\end{itemize}
继续组合 $S$ 和 $T$，轨道中有无穷多个点，密密麻麻地分布在上半平面中。
这些点在 $\SL_2(\Z)$ 的视角下全是"同一个点"。

基本域的作用就是：从这一大堆等价的点里，选出一个"标准代表"。
对 $2i$ 来说，它自己就在基本域中（$|\Re(2i)| = 0 < 1/2$，$|2i| = 2 > 1$），
所以 $2i$ 就是这个轨道的标准代表。
\end{example}

\begin{definition}[基本域]
\label{def:fundamental-domain}
\index{基本域 (fundamental domain)}
群 $G$ 作用在空间 $X$ 上时，一个\textbf{基本域}（fundamental domain）是
$X$ 的子集 $\mathcal{F}$，使得：
\begin{enumerate}
  \item $X = \bigcup_{\gamma \in G} \gamma(\overline{\mathcal{F}})$
    （覆盖性）；
  \item $\gamma(\mathcal{F}) \cap \mathcal{F} = \emptyset$ 对所有
    $\gamma \neq \Id$（内部不重叠）。
\end{enumerate}
\end{definition}

\begin{theorem}[$\SL_2(\Z)$ 的标准基本域]
\label{thm:fundamental-domain}
\[
  \mathcal{F} = \left\{\tau \in \HH \;\middle|\;
    |\tau| > 1,\; -\frac{1}{2} < \Re(\tau) < \frac{1}{2}
  \right\}
\]
（即开域，边界按约定选取一半）是 $\SL_2(\Z)$ 作用在 $\HH$ 上的基本域。
更精确地：$\overline{\mathcal{F}}$ 的闭包为
\[
  \overline{\mathcal{F}} = \left\{\tau \in \HH \;\middle|\;
    |\tau| \geq 1,\; -\frac{1}{2} \leq \Re(\tau) \leq \frac{1}{2}
  \right\},
\]
且 $\HH = \bigcup_{\gamma \in \SL_2(\Z)} \gamma(\overline{\mathcal{F}})$。
\end{theorem}

\begin{proof}
\textbf{覆盖性}：给定 $\tau \in \HH$，我们找 $\gamma \in \SL_2(\Z)$
使得 $\gamma(\tau) \in \overline{\mathcal{F}}$。

由\cref{prop:mobius-preserves-H}，$\im(\gamma(\tau)) = \im(\tau)/|c\tau + d|^2$。
对固定的 $\tau$，$|c\tau + d|^2 = (cx + d)^2 + c^2 y^2 \geq c^2 y^2$，
只有有限多对 $(c,d)$ 使得 $|c\tau + d| \leq 1/y$（即使虚部增大）。
因此 $\{\im(\gamma(\tau)) \mid \gamma \in \SL_2(\Z)\}$ 有最大值，
设在 $\gamma_0$ 处取到。

先用 $T^n$（$n$ 适当）使得 $\Re(\gamma_0(\tau)) \in [-1/2, 1/2]$。
设结果为 $\tau' = T^n \gamma_0(\tau)$。

若 $|\tau'| < 1$，则 $S(\tau') = -1/\tau'$ 满足
$\im(S(\tau')) = \im(\tau')/|\tau'|^2 > \im(\tau')$，
矛盾于 $\im(\tau')$ 已是最大值。故 $|\tau'| \geq 1$，
即 $\tau' \in \overline{\mathcal{F}}$。

\textbf{内部不重叠}：设 $\tau, \tau' \in \mathcal{F}$（开域内部），
$\gamma(\tau) = \tau'$，$\gamma = \begin{pmatrix} a & b \\ c & d \end{pmatrix}
\neq \pm I$。不妨设 $\im(\tau') \geq \im(\tau)$，即
$|c\tau + d|^2 \leq 1$。

由 $\tau \in \mathcal{F}$，$\im(\tau) > \sqrt{3}/2$。
若 $|c| \geq 2$，则 $|c\tau + d| \geq |c|\im(\tau) > \sqrt{3} > 1$，矛盾。
若 $|c| = 1$，则 $|c\tau + d|^2 = |\tau + d|^2 \leq 1$（当 $c = 1$），
但 $|\tau + d|^2 = (x + d)^2 + y^2$。由 $|x| < 1/2$ 和 $y > \sqrt{3}/2$，
\[
  |\tau + d|^2 \geq y^2 > 3/4.
\]
若 $d = 0$，则 $|\tau|^2 \leq 1$，但 $\tau \in \mathcal{F}$ 要求
$|\tau| > 1$，矛盾。
类似分析 $d = \pm 1$ 的情形均导致 $\tau$ 在 $\mathcal{F}$ 的边界上而非内部。
若 $c = 0$，则 $\gamma = \pm T^b$，$\tau' = \tau \pm b$，
由 $|\Re(\tau)|, |\Re(\tau')| < 1/2$ 得 $b = 0$，$\gamma = \pm I$。
\end{proof}

\textbf{基本域的形状。}
$\overline{\mathcal{F}}$ 形如一个无限高的"拱门"：底部弧线是单位圆
$|\tau| = 1$ 的一段，两侧是 $\Re(\tau) = \pm 1/2$ 的垂线。
这个形状正好对应 $S$ 和 $T$ 的分界线：$T$ 把竖直带平移，
所以左右边界由 $T$ 粘合；$S$ 关于单位圆反演，所以底部弧线由 $S$ 粘合。

\begin{example}[判断一个点是否在基本域中]
\label{ex:point-in-F}
\leavevmode
\begin{enumerate}
  \item $\tau_1 = \frac{1}{5} + 2i$：
    $|\Re(\tau_1)| = 1/5 < 1/2$，$|\tau_1|^2 = 1/25 + 4 > 1$。
    两个条件都满足，所以 $\tau_1 \in \overline{\mathcal{F}}$。

  \item $\tau_2 = 3 + i$：
    $|\Re(\tau_2)| = 3 > 1/2$，不在基本域中。
    但 $T^{-3}(\tau_2) = \tau_2 - 3 = i$，而 $i$ 在 $\overline{\mathcal{F}}$
    的边界上（$|i| = 1$）。所以 $\tau_2$ 和 $i$ 属于同一个轨道。

  \item $\tau_3 = i/3$：
    $|\tau_3| = 1/3 < 1$，不在基本域中。
    用 $S$：$S(\tau_3) = -1/(i/3) = -3/i = 3i$。
    而 $3i$ 满足 $|\Re(3i)| = 0 < 1/2$，$|3i| = 3 > 1$，
    所以 $3i \in \mathcal{F}$。
\end{enumerate}
\end{example}

\begin{example}[把一个点"折回"基本域]
\label{ex:reduce-to-F}
取 $\tau = \frac{7}{3} + \frac{i}{2}$。怎样找到它在基本域中的代表元？

\textbf{第一步}：用 $T$ 调整实部。$\Re(\tau) = 7/3 \approx 2.33$，
取 $T^{-2}$：$\tau' = \tau - 2 = \frac{1}{3} + \frac{i}{2}$。
现在 $|\Re(\tau')| = 1/3 < 1/2$，实部条件满足了。

\textbf{第二步}：检查模长。$|\tau'|^2 = 1/9 + 1/4 = 13/36 < 1$，
不满足 $|\tau| \geq 1$。用 $S$：
\[
  S(\tau') = \frac{-1}{\frac{1}{3} + \frac{i}{2}}
  = \frac{-1}{\frac{2 + 3i}{6}} = \frac{-6}{2 + 3i}
  = \frac{-6(2 - 3i)}{4 + 9} = \frac{-12 + 18i}{13}.
\]
所以 $S(\tau') = -\frac{12}{13} + \frac{18}{13}i$。

\textbf{第三步}：再用 $T$ 调整实部。$\Re(S(\tau')) = -12/13 \approx -0.92$，
取 $T^1$：$\tau'' = S(\tau') + 1 = \frac{1}{13} + \frac{18}{13}i$。
现在 $|\Re(\tau'')| = 1/13 < 1/2$，
$|\tau''|^2 = 1/169 + 324/169 = 325/169 = 25/13 > 1$。两个条件都满足！

结论：$\tau'' = \frac{1}{13} + \frac{18}{13}i \in \mathcal{F}$，
由 $\tau'' = TS T^{-2}(\tau)$ 得到。
\end{example}

\textbf{椭圆不动点与稳定子。}
\label{rem:elliptic-fixed-points}
大多数 $\tau \in \HH$ 只有"平凡"的对称性——只有 $\pm I$
把它映回自身。但基本域边界上有两个特殊点，它们有额外的对称性。

给定 $\tau$，所有把 $\tau$ 固定不动的矩阵构成一个子群，
叫做 $\tau$ 的\textbf{稳定子}\index{稳定子 (stabilizer)}（stabilizer）：
$(\SL_2(\Z))_\tau = \{\gamma \in \SL_2(\Z) \mid \gamma(\tau) = \tau\}$。

\begin{enumerate}
  \item $\tau = i$：$S(i) = -1/i = i$，所以 $S$ 固定 $i$。
    稳定子为 $\{I, S, -I, -S\}$，同构于 $\Z/4\Z$
    （在 $\PSL_2(\Z)$ 中为 $\Z/2\Z$）。
  \item $\tau = \rho = e^{2\pi i/3} = \frac{-1+i\sqrt{3}}{2}$：
    直接计算 $(ST)(\rho) = S(\rho + 1) = \rho$。
    稳定子为 $\langle ST \rangle \cong \Z/6\Z$
    （在 $\PSL_2(\Z)$ 中为 $\Z/3\Z$）。
\end{enumerate}
这两个\textbf{椭圆不动点}导致基本域边界上的粘合出现特殊行为，
在\cref{ch:dimension}的维数公式中会产生修正项。


% ============================================================================

\section{尖点}
\label{sec:cusps}
% ============================================================================

基本域 $\mathcal{F}$ 向上延伸到无穷远——它的"顶部"越来越窄，
像一个无限高的烟囱。这个"烟囱口"在数学上需要一个名字，
它就是\textbf{尖点}（cusp）。

\textbf{为什么需要尖点？}
回忆模形式 $E_4(\tau)$ 的 $q$-展开 $\sum a_n q^n$，其中 $q = e^{2\pi i\tau}$。
当 $\im(\tau) \to +\infty$ 时，$q \to 0$。
所以"$\tau$ 趋向无穷远"对应"$q$ 趋向 $0$"，而我们要求模形式在 $q = 0$
处有良好的行为（全纯，甚至为零）。
这个"无穷远点"就是最基本的尖点 $\infty$。

\textbf{$\mathbb{P}^1(\Q)$ 是什么？}
除了 $\infty$ 之外，还有其他"方向"可以趋向实轴上的有理点。
$\mathbb{P}^1(\Q)$ 就是 $\Q \cup \{\infty\}$——
所有有理数加上一个无穷远点，可以理解为"有理数的圆周"。

\begin{definition}[扩展上半平面与尖点]
\label{def:cusps}
\index{扩展上半平面}\index{尖点 (cusp)}\index{P1Q@$\mathbb{P}^1(\mathbb{Q})$}
\textbf{扩展上半平面}定义为
\[
  \HH^* = \HH \cup \Q \cup \{\infty\} = \HH \cup \mathbb{P}^1(\Q).
\]
$\SL_2(\Z)$ 的莫比乌斯变换自然延拓到 $\mathbb{P}^1(\Q)$ 上：
对 $\gamma = \begin{pmatrix} a & b \\ c & d \end{pmatrix}$，
\[
  \gamma(\infty) = \frac{a}{c}, \qquad
  \gamma\!\left(-\frac{d}{c}\right) = \infty.
\]
（就是把 $\frac{a\tau + b}{c\tau + d}$ 中分母为零和分子为零的情况补上。）

$\mathbb{P}^1(\Q)$ 中的点称为\textbf{尖点}（cusps）。
\end{definition}

\begin{example}[具体的尖点变换]
\label{ex:cusp-action}
\leavevmode
\begin{itemize}
  \item $T(\infty) = \infty + 1 = \infty$
    ——平移不影响无穷远点。
  \item $S(\infty) = -1/\infty = 0$
    ——反演把 $\infty$ 映到 $0$。
  \item $S(0) = -1/0 = \infty$
    ——反演把 $0$ 映回 $\infty$。
  \item 取 $\gamma = \begin{pmatrix} 3 & 1 \\ 2 & 1 \end{pmatrix}
    \in \SL_2(\Z)$（$3 \cdot 1 - 1 \cdot 2 = 1$），
    则 $\gamma(\infty) = 3/2$。
    所以 $3/2$ 和 $\infty$ 在同一轨道中。
\end{itemize}
\end{example}

\begin{proposition}
\label{prop:cusps-transitive}
$\SL_2(\Z)$ 在 $\mathbb{P}^1(\Q)$ 上的作用是传递的，即所有尖点都等价。
换句话说：$\SL_2(\Z) \backslash \mathbb{P}^1(\Q)$ 只有一个点。
\end{proposition}

\begin{proof}
任取 $a/c \in \Q$（$\gcd(a,c) = 1$），由 Bézout 恒等式，存在 $b, d \in \Z$
使得 $ad - bc = 1$。则
$\gamma = \begin{pmatrix} a & b \\ c & d \end{pmatrix} \in \SL_2(\Z)$
满足 $\gamma(\infty) = a/c$。所以每个有理数都和 $\infty$ 在同一轨道中。
\end{proof}

\textbf{同余子群的尖点。}
对 $\SL_2(\Z)$ 来说只有一个尖点等价类，但对更小的同余子群
（如 $\GammaO(N)$），尖点会分裂成多个等价类。
例如 $\GammaO(2)$ 有 $2$ 个不等价的尖点：$\infty$ 和 $0$
（因为 $S \notin \GammaO(2)$，无法把 $\infty$ 映到 $0$）。
这将在\cref{sec:congruence-subgroups}中进一步讨论。


% ============================================================================

\section{模形式的定义}
\label{sec:modular-forms-def}
% ============================================================================

在章开头，我们已经看到 $E_4(\tau)$ 满足 $E_4(\tau+1) = E_4(\tau)$ 和
$E_4(-1/\tau) = \tau^4 E_4(\tau)$。现在我们把这种对称性抽象为一般定义。

\textbf{只有变换律还不够——为什么需要额外的"尖点条件"？}

仅满足变换律 $f(\gamma\tau) = (c\tau+d)^k f(\tau)$ 的全纯函数，
我们称之为"弱模形式"。但如果不加其他限制，这样的函数太多了——
比如 $1/\Delta(\tau)$ 在尖点处有一阶极点，它满足变换律，
但代数性质很差，不适合用来做算术。

\textbf{$q$-展开与尖点全纯条件。}
由 $f(\tau+1) = f(\tau)$，$f$ 可以写成 $q = e^{2\pi i\tau}$ 的 Fourier 级数：
\[
  f(\tau) = \sum_{n=-\infty}^{+\infty} a_n q^n.
\]
当 $\im(\tau) \to +\infty$ 时，$q \to 0$。
"$f$ 在尖点全纯"等价于 $f$ 作为 $q$ 的函数在 $q=0$ 处没有极点，
也就是 $a_n = 0$ 对所有 $n < 0$。

\textbf{这个条件的效果：}它把无限维的函数空间压缩成有限维！
一个满足变换律且在尖点全纯的函数必须同时满足太多约束，
以至于这样的函数空间只有有限维（我们在\cref{ch:dimension}中会证明这个维数公式）。
有限维才能用线性代数的工具，才能谈"基"和"特征值"。

\textbf{为什么尖点形式更特殊？}
尖点形式额外要求 $a_0 = 0$，即 $f$ 在无穷远处"消失"。
这让尖点形式具有最好的分析性质（$f(\tau) = O(e^{-2\pi y})$ 当 $y\to\infty$），
特别适合定义 Petersson 内积（\cref{sec:petersson}）和研究 $L$-函数。
重量 $12$ 的判别形式 $\Delta$ 就是尖点形式的典型代表。

\begin{definition}[权 $k$ 的弱模形式]
\label{def:weakly-modular}
\index{弱模形式 (weakly modular form)}
设 $k \in \Z$。全纯函数 $f\colon \HH \to \C$ 称为权 $k$ 的
\textbf{弱模形式}（weakly modular form），若对所有
$\gamma = \begin{pmatrix} a & b \\ c & d \end{pmatrix} \in \SL_2(\Z)$，
\[
  f(\gamma(\tau)) = (c\tau + d)^k f(\tau).
\]
由\cref{thm:sl2z-generators}，这等价于同时满足：
\begin{enumerate}
  \item $f(\tau + 1) = f(\tau)$（由 $T$ 给出）；
  \item $f(-1/\tau) = \tau^k f(\tau)$（由 $S$ 给出）。
\end{enumerate}
\end{definition}

条件 (1) 意味着 $f$ 有以 $q = e^{2\pi i\tau}$ 为变量的 Fourier 展开，
称为 \textbf{$q$-展开}\index{q-展开 ($q$-expansion)}（$q$-expansion）：
\[
  f(\tau) = \sum_{n \in \Z} a_n q^n, \quad q = e^{2\pi i\tau}.
\]

\begin{definition}[模形式与尖点形式]
\label{def:modular-form}
\index{模形式 (modular form)}\index{尖点形式 (cusp form)}\index{Mk@$\mathcal{M}_k$}\index{Sk@$\mathcal{S}_k$}
权 $k$ 的\textbf{模形式}（modular form）是权 $k$ 的弱模形式 $f$，
且 $f$ 在尖点 $\infty$ 处\textbf{全纯}，即 $q$-展开中 $a_n = 0$ 对所有 $n < 0$：
\[
  f(\tau) = \sum_{n=0}^{\infty} a_n q^n.
\]
权 $k$ 的模形式空间记为 $\Mk_k(\SL_2(\Z))$。

若进一步 $a_0 = 0$（即 $f$ 在尖点处的值为零），则 $f$ 称为
\textbf{尖点形式}（cusp form）。尖点形式空间记为 $\Sk_k(\SL_2(\Z))$。
\end{definition}

\begin{example}[Eisenstein 级数]
\label{ex:eisenstein-preview}
对偶整数 $k \geq 4$，\textbf{Eisenstein 级数}
\[
  G_k(\tau) = \sum_{\substack{(c,d) \in \Z^2 \\ (c,d) \neq (0,0)}}
  \frac{1}{(c\tau + d)^k}
\]
是权 $k$ 的模形式。我们将在\cref{ch:eisenstein}中详细讨论。
规范化版本 $E_k = G_k / (2\zeta(k))$ 的 $q$-展开为
\[
  E_k(\tau) = 1 - \frac{2k}{B_k} \sum_{n=1}^{\infty} \sigma_{k-1}(n) q^n,
\]
其中 $B_k$ 是 Bernoulli 数，$\sigma_{k-1}(n) = \sum_{d \mid n} d^{k-1}$。
\end{example}

\begin{example}[判别形式 $\Delta$]
\label{ex:discriminant-form}
\textbf{判别形式}
\[
  \Delta(\tau) = \frac{1}{1728}\bigl(E_4(\tau)^3 - E_6(\tau)^2\bigr)
  = q \prod_{n=1}^{\infty}(1 - q^n)^{24}
  = \sum_{n=1}^{\infty} \tau(n) q^n
\]
是权 $12$ 的尖点形式，其中 $\tau(n)$ 是 \textbf{Ramanujan $\tau$-函数}。
$\Delta$ 在 $\HH$ 上无零点。
\end{example}

\begin{example}[$j$-不变量]
\label{ex:j-invariant-function}
\textbf{$j$-不变量}
\[
  j(\tau) = \frac{E_4(\tau)^3}{\Delta(\tau)} = \frac{1}{q} + 744 + 196884q + \cdots
\]
是 $\SL_2(\Z)$ 的一个\textbf{模函数}（权 $0$ 的亚纯模形式），
它建立了 $\SL_2(\Z) \backslash \HH$ 到 $\C$ 的双射。
\end{example}


% ============================================================================

\section{同余子群}
\label{sec:congruence-subgroups}
% ============================================================================

\textbf{我们为什么需要同余子群？先看一个具体的算术问题。}

有一条著名的椭圆曲线：
\[
  E\colon y^2 = x^3 - x.
\]
它有一个特殊的有理 $4$ 阶点，$P = (0, 0)$（验证：$2P = (0,0)$ 到 ... 事实上 $P$ 是 $2$ 阶点，
我们换一个例子）。

更深刻的问题是：能不能用模形式来捕捉"椭圆曲线 $E$ 带着某个 $N$ 阶子群 $C$"这个结构？

用 $\SL_2(\Z)$ 对应的模形式空间 $\Mk_k(\SL_2(\Z))$，
我们只能看到椭圆曲线的同构类——不区分有没有额外结构。
如果想区分"有 $N$ 阶子群 $C$"和"没有 $N$ 阶子群"的椭圆曲线，
就需要用更小的对称群，只保留"不破坏 $C$ 的结构"的变换。

这正是\textbf{同余子群}（congruence subgroup）的来源：
$\Gamma_0(N)$、$\Gamma_1(N)$、$\Gamma(N)$ 是 $\SL_2(\Z)$ 的子群，
它们对应"记住 $N$ 阶额外结构"的那部分对称性。

\textbf{直觉上，三种子群对应三种额外信息：}
\begin{itemize}
  \item $\Gamma_0(N)$：记住一个 $N$ 阶循环子群 $C \subset E[N]$（但不区分 $C$ 中的元素）。
    条件是矩阵左下角 $c \equiv 0 \pmod N$。
  \item $\Gamma_1(N)$：记住 $C$ 中一个具体的 $N$ 阶点 $P$（不只是子群，而是点）。
    条件是 $a \equiv d \equiv 1$，$c \equiv 0 \pmod N$。
  \item $\Gamma(N)$：记住 $E[N]$（$E$ 的所有 $N$ 阶点）的一整组有序基 $(P, Q)$。
    条件是整个矩阵 $\equiv I \pmod N$。
\end{itemize}
越小的子群记住的信息越多，对应的模形式空间也就越大。
这将在\cref{thm:moduli-interpretation}中精确表述。

\textbf{同余条件"$\pmod N$"从哪来？}
直觉：如果我们只关心 $E$ 的 $N$ 阶点，
那么在格 $\Lambda = \Z\tau + \Z$ 上，$N$ 阶点就是 $\frac{1}{N}\Lambda / \Lambda$。
一个矩阵 $\gamma = \begin{pmatrix}a&b\\c&d\end{pmatrix}$ 作用在这些点上，
它的效果只取决于 $\gamma \bmod N$——所以"记住 $N$ 阶点结构"对应的子群
自然就是"在模 $N$ 意义下满足某些条件的矩阵"。
\begin{definition}[同余子群]
\label{def:congruence-subgroups}
\index{同余子群 (congruence subgroup)}\index{Gamma0@$\Gamma_0(N)$}\index{Gamma1@$\Gamma_1(N)$}\index{GammaN@$\Gamma(N)$}
设 $N$ 为正整数。定义以下 $\SL_2(\Z)$ 的子群：
\begin{align}
  \Gamma(N) &= \left\{
    \begin{pmatrix} a & b \\ c & d \end{pmatrix} \in \SL_2(\Z)
    \;\middle|\;
    \begin{pmatrix} a & b \\ c & d \end{pmatrix}
    \equiv \begin{pmatrix} 1 & 0 \\ 0 & 1 \end{pmatrix} \pmod{N}
  \right\}, \\[6pt]
  \GammaI(N) &= \left\{
    \begin{pmatrix} a & b \\ c & d \end{pmatrix} \in \SL_2(\Z)
    \;\middle|\;
    a \equiv d \equiv 1,\; c \equiv 0 \pmod{N}
  \right\}, \\[6pt]
  \GammaO(N) &= \left\{
    \begin{pmatrix} a & b \\ c & d \end{pmatrix} \in \SL_2(\Z)
    \;\middle|\;
    c \equiv 0 \pmod{N}
  \right\}.
\end{align}
包含关系为 $\Gamma(N) \trianglelefteq \GammaI(N) \trianglelefteq \GammaO(N)
\leq \SL_2(\Z)$。

更一般地，$\SL_2(\Z)$ 的子群 $\Gamma$ 若包含某个 $\Gamma(N)$，
则称为\textbf{同余子群}（congruence subgroup），满足此条件的最小 $N$ 称为
$\Gamma$ 的\textbf{级}（level）。
\end{definition}

\begin{proposition}[指标公式]
\label{prop:index-formulas}
\leavevmode
\begin{enumerate}
  \item $[\SL_2(\Z) : \Gamma(N)] = N^3 \prod_{p \mid N}(1 - 1/p^2)$
    （$N \geq 1$）。
  \item $[\SL_2(\Z) : \GammaI(N)] = N^2 \prod_{p \mid N}(1 - 1/p^2)$
    （$N \geq 1$）。
  \item $[\SL_2(\Z) : \GammaO(N)] = N \prod_{p \mid N}(1 + 1/p)$
    （$N \geq 1$）。
\end{enumerate}
\end{proposition}

\begin{proof}
考虑约化同态 $\SL_2(\Z) \to \SL_2(\Z/N\Z)$。

\textbf{(1)}：$\Gamma(N) = \ker(\SL_2(\Z) \to \SL_2(\Z/N\Z))$。
此约化同态是满射（这需要证明，我们在习题中给出提示）。
故 $[\SL_2(\Z) : \Gamma(N)] = \#\SL_2(\Z/N\Z)$。
由中国剩余定理，$\#\SL_2(\Z/N\Z) = \prod_{p^e \| N} \#\SL_2(\Z/p^e\Z)$。
而 $\#\SL_2(\Z/p^e\Z) = p^{3e} \prod_{p \mid p^e}(1 - 1/p^2)
= p^{3(e-1)} \cdot p^3(1 - 1/p^2) = p^{3e}(1 - 1/p^2)$。

\textbf{(2)}：$\GammaI(N)/\Gamma(N) \cong \Z/N\Z$
（通过映射 $\gamma \mapsto b \bmod N$），故
$[\SL_2(\Z) : \GammaI(N)] = [\SL_2(\Z) : \Gamma(N)]/N$。

\textbf{(3)}：$\GammaO(N)/\GammaI(N) \cong (\Z/N\Z)^\times$
（通过映射 $\gamma \mapsto d \bmod N$），故
$[\SL_2(\Z) : \GammaO(N)] = [\SL_2(\Z) : \GammaI(N)]/\varphi(N)$。
化简得 $N \prod_{p \mid N}(1 + 1/p)$。
\end{proof}

\begin{definition}[同余子群上的模形式]
\label{def:modular-forms-congruence}
设 $\Gamma$ 为同余子群。权 $k$ 的\textbf{模形式} $f \in \Mk_k(\Gamma)$
是全纯函数 $f\colon \HH \to \C$，满足：
\begin{enumerate}
  \item 对所有 $\gamma = \begin{pmatrix} a & b \\ c & d \end{pmatrix}
    \in \Gamma$，
    \[
      f(\gamma(\tau)) = (c\tau + d)^k f(\tau);
    \]
  \item $f$ 在 $\Gamma$ 的每个尖点处全纯。
\end{enumerate}
若 $f$ 在所有尖点处的值为零（即每个尖点的 $q$-展开常数项 $a_0 = 0$），则 $f$ 是\textbf{尖点形式}，$f \in \Sk_k(\Gamma)$。
\end{definition}


% ============================================================================

\section{复环面}
\label{sec:complex-tori}
% ============================================================================

前面我们一直在研究上半平面 $\HH$ 上的函数（模形式）
以及 $\SL_2(\Z)$ 在 $\HH$ 上的作用（基本域、尖点）。
现在要回答一个自然的问题：
$\SL_2(\Z) \backslash \HH$（所有轨道的集合）到底"是"什么空间？

答案是：它参数化了所有的\textbf{复环面}。
复环面是最简单的、非平凡的紧黎曼面，也是椭圆曲线的复分析版本。

\textbf{什么是格？}
在复平面 $\C$ 中，选两个 $\R$-线性无关的复数 $\omega_1, \omega_2$
（即 $\omega_1/\omega_2 \notin \R$），它们生成一个"网格"：
\[
  \Lambda = \Z\omega_1 + \Z\omega_2
  = \{m\omega_1 + n\omega_2 \mid m, n \in \Z\}.
\]

\begin{example}[几个具体的格]
\label{ex:lattices}
\leavevmode
\begin{itemize}
  \item $\Lambda = \Z + \Z i$：正方形格，格点分布在整数坐标上。
    基本平行四边形是单位正方形。
  \item $\Lambda = \Z + \Z e^{2\pi i/3}$：正三角形格（蜂巢状），
    $\omega_2 = -1/2 + i\sqrt{3}/2$。
    基本平行四边形是 $60°$ 的菱形。
  \item $\Lambda = \Z + \Z(2i)$：长方形格，
    基本平行四边形是 $1 \times 2$ 的长方形。
\end{itemize}
\end{example}

\begin{definition}[格与复环面]
\label{def:lattice-torus}
\index{格 (lattice)}\index{复环面 (complex torus)}
$\C$ 中的一个\textbf{格}（lattice）是
$\Lambda = \Z\omega_1 + \Z\omega_2$（$\omega_1/\omega_2 \notin \R$）。

对应的\textbf{复环面}为商群
\[
  \C/\Lambda = \{z + \Lambda \mid z \in \C\}.
\]
两个复数 $z_1, z_2$ 在 $\C/\Lambda$ 中代表同一个点，
当且仅当 $z_1 - z_2 \in \Lambda$（即它们差一个格向量）。
\end{definition}

\textbf{复环面长什么样？}
可以这样想象：取格的一个\textbf{基本平行四边形}
（以 $0, \omega_1, \omega_2, \omega_1+\omega_2$ 为顶点），
然后把对边粘合——上下粘合得到一个圆柱，再把圆柱的两端粘合得到一个甜甜圈（torus）。
这就是为什么叫"环面"。

拓扑上，$\C/\Lambda$ 是一个亏格为 $1$ 的紧曲面。
但它不仅有拓扑结构——从 $\C$ 继承的复结构使它成为一个\textbf{黎曼面}
（Riemann surface，详见\cref{ch:riemann-surfaces}），
而且加法 $z + w$ 使它成为一个\textbf{复 Lie 群}（即加法群）。

\begin{example}[格 $\Z + \Z i$ 对应的环面]
\label{ex:square-torus}
取 $\Lambda = \Z + \Z i$。基本平行四边形是单位正方形
$\{x + iy \mid 0 \leq x < 1,\; 0 \leq y < 1\}$。
\begin{itemize}
  \item 点 $0.3 + 0.7i$ 和 $1.3 + 0.7i$ 在 $\C/\Lambda$ 中是同一个点
    （差 $1 \in \Lambda$）。
  \item 点 $0.5 + 0.5i$ 和 $-0.5 - 0.5i$ 也是同一个点
    （差 $1 + i \in \Lambda$）。
  \item 加法：$(0.6 + 0.8i) + (0.7 + 0.5i) = 1.3 + 1.3i \equiv 0.3 + 0.3i$
    （模掉 $1 + i$）。
\end{itemize}
\end{example}

\begin{proposition}[复环面的同构分类]
\label{prop:torus-isomorphism}
$\C/\Lambda \cong \C/\Lambda'$（作为复 Lie 群）当且仅当
$\Lambda' = \alpha\Lambda$ 对某 $\alpha \in \C^*$。
\end{proposition}

\begin{proof}
设 $\varphi\colon \C/\Lambda \to \C/\Lambda'$ 为全纯群同构。

\textbf{关键思路：从"卷起来的"回到"展开的"。}
复环面 $\C/\Lambda$ 是把复平面 $\C$ 按照格 $\Lambda$ "卷起来"得到的
（就像把一张纸沿两个方向卷成甜甜圈）。
反过来，$\C$ 就是 $\C/\Lambda$ "展开后"的样子——
拓扑学中叫做\textbf{万有覆盖}（universal cover）。
投影映射 $\pi\colon \C \to \C/\Lambda$（$z \mapsto z + \Lambda$）
把 $\C$ 中相差一个格向量 $\omega \in \Lambda$ 的点映到同一个点。

现在，$\varphi$ 是环面之间的映射。我们想知道它"展开后"是什么样的。
覆盖空间理论保证：$\varphi$ 可以"提升"为一个全纯函数
$\tilde\varphi\colon \C \to \C$，使得下图交换：
\[
  \begin{tikzcd}
    \C \arrow[r, "\tilde\varphi"] \arrow[d, "\pi"'] & \C \arrow[d, "\pi'"] \\
    \C/\Lambda \arrow[r, "\varphi"'] & \C/\Lambda'
  \end{tikzcd}
\]
即 $\pi' \circ \tilde\varphi = \varphi \circ \pi$。
直觉上：在展开的平面上做的事情，卷回去之后要和 $\varphi$ 一致。

\textbf{推导 $\tilde\varphi$ 的形状。}
由于 $\pi(z + \omega) = \pi(z)$（差一个格向量不影响），
$\tilde\varphi(z + \omega)$ 和 $\tilde\varphi(z)$ 在 $\C/\Lambda'$
中代表同一个点，即
$\tilde\varphi(z + \omega) - \tilde\varphi(z) \in \Lambda'$。

函数 $z \mapsto \tilde\varphi(z + \omega) - \tilde\varphi(z)$ 是连续的，
但它的值只能落在离散集 $\Lambda'$ 中——
连续函数的像如果是离散的，就必须是常数。
所以对每个固定的 $\omega$，这个差是一个不依赖于 $z$ 的常数。

对 $z$ 求导：$\tilde\varphi'(z + \omega) = \tilde\varphi'(z)$，
即 $\tilde\varphi'$ 以 $\Lambda$ 为周期。
一个全纯的双周期函数（椭圆函数）如果没有极点，
由 Liouville 定理必为常数（\cref{thm:elliptic-basic}(1)）。
所以 $\tilde\varphi'(z) = \alpha$ 是常数，积分得 $\tilde\varphi(z) = \alpha z + \beta$。

最后，$\tilde\varphi$ 把 $\Lambda$ 映到 $\Lambda'$ 内（因为差值在 $\Lambda'$ 中），
又由满射性得 $\alpha\Lambda = \Lambda'$。
\end{proof}

\begin{corollary}
\label{cor:torus-tau}
每个复环面同构于 $\C/\Lambda_\tau$（$\Lambda_\tau = \Z\tau + \Z$，
$\tau \in \HH$）。两个环面 $\C/\Lambda_\tau$ 和 $\C/\Lambda_{\tau'}$
同构当且仅当 $\tau' = \gamma(\tau)$ 对某 $\gamma \in \SL_2(\Z)$。

因此，\textbf{复环面的同构类由 $\SL_2(\Z) \backslash \HH$ 的点参数化}。
这就是上半平面与模群的几何意义——它们在"管理"所有复环面。
\end{corollary}

\begin{proof}
\textbf{第一步：标准化格。}
任意格 $\Lambda = \Z\omega_1 + \Z\omega_2$。
不妨设 $\im(\omega_1/\omega_2) > 0$（否则交换 $\omega_1, \omega_2$）。

乘以 $\alpha = 1/\omega_2$：$\alpha\Lambda = \Z(\omega_1/\omega_2) + \Z$。
令 $\tau = \omega_1/\omega_2$，则 $\alpha\Lambda = \Z\tau + \Z = \Lambda_\tau$，
且 $\tau \in \HH$。
由\cref{prop:torus-isomorphism}，$\C/\Lambda \cong \C/\Lambda_\tau$。

\textbf{关键直觉}：任何格都可以通过缩放变成"第一个生成元在上半平面、
第二个生成元为 $1$"的标准形式。

\textbf{第二步：什么时候两个标准格相同？}
$\C/\Lambda_\tau \cong \C/\Lambda_{\tau'}$ 当且仅当
$\Lambda_{\tau'} = \alpha\Lambda_\tau$ 对某 $\alpha \in \C^*$。

$\Lambda_\tau = \Z\tau + \Z$，所以 $\alpha\Lambda_\tau = \Z(\alpha\tau) + \Z\alpha$。
要求这等于 $\Lambda_{\tau'} = \Z\tau' + \Z$，
即 $\alpha\tau$ 和 $\alpha$ 都是 $\tau'$ 和 $1$ 的整系数线性组合：
\[
  \alpha\tau = a\tau' + b, \qquad \alpha = c\tau' + d,
  \qquad a,b,c,d \in \Z.
\]
消去 $\alpha$：$\tau = \frac{a\tau' + b}{c\tau' + d}$。
由 $\im(\tau) > 0$、$\im(\tau') > 0$ 推出 $ad - bc > 0$。
反过来 $\tau'$ 也要能用 $\tau$ 表示（因为是同构不是单方向），
所以 $ad - bc = 1$，即 $\begin{pmatrix} a & b \\ c & d \end{pmatrix} \in \SL_2(\Z)$。
\end{proof}

\begin{example}[同一个环面的不同参数]
\label{ex:equivalent-tau}
取 $\tau = 2i$ 和 $\tau' = -1/(2i) = i/2$。
它们满足 $\tau' = S(\tau)$（$S \in \SL_2(\Z)$），
所以 $\C/\Lambda_{2i} \cong \C/\Lambda_{i/2}$。
具体地：$\Lambda_{2i} = \Z(2i) + \Z$，
取 $\alpha = 1/(2i) = -i/2$，则
$\alpha \cdot \Lambda_{2i} = \Z(1) + \Z(-i/2) = \Z + \Z(i/2) = \Lambda_{i/2}$。
\end{example}


% ============================================================================

\section{复环面与椭圆曲线}
\label{sec:tori-and-ec}
% ============================================================================

复环面 $\C/\Lambda$ 是一个"分析对象"——用复平面和格定义的。
\textbf{椭圆曲线}是一个"代数对象"——用多项式方程定义的。
本节的核心结论是：这两种看似不同的东西其实是同一回事。

\textbf{什么是椭圆曲线？}
最简单地说，椭圆曲线就是形如
\[
  E\colon y^2 = x^3 + ax + b \qquad (\text{其中 } 4a^3 + 27b^2 \neq 0)
\]
的代数曲线，加上一个"无穷远点" $O$。
条件 $4a^3 + 27b^2 \neq 0$ 保证曲线没有尖点或自交
（即它是"光滑"的）。
椭圆曲线上的点可以定义一个加法群结构——
这是代数几何中最神奇的事情之一，但这里我们不展开。

\textbf{连接的桥梁：Weierstrass $\wp$ 函数。}
给定格 $\Lambda$，\textbf{Weierstrass $\wp$ 函数}定义为
\[
  \wp(z;\Lambda) = \frac{1}{z^2} + \sum_{\omega \in \Lambda \setminus \{0\}}
  \left(\frac{1}{(z-\omega)^2} - \frac{1}{\omega^2}\right).
\]
它是一个以 $\Lambda$ 为周期的亚纯函数（椭圆函数），
在每个格点 $\omega \in \Lambda$ 处有二阶极点，其他地方全纯。
关键性质是 $\wp$ 满足微分方程：
\[
  (\wp')^2 = 4\wp^3 - g_2\wp - g_3,
\]
其中 $g_2 = 60\sum_{\omega \neq 0} \omega^{-4}$，
$g_3 = 140\sum_{\omega \neq 0} \omega^{-6}$。
这个微分方程正好是一条椭圆曲线！

\begin{theorem}[均匀化定理]
\label{thm:uniformization}
\leavevmode
\begin{enumerate}
  \item \textbf{从环面到曲线}：映射
    \[
      \C/\Lambda \to E_\Lambda(\C), \quad
      z \mapsto (\wp(z), \wp'(z))
    \]
    （$z = 0$ 映到无穷远点 $O$）是一个群同构。
    这里 $E_\Lambda\colon y^2 = 4x^3 - g_2 x - g_3$。
    
    直觉：$\wp$ 和 $\wp'$ 把环面上的每个点
    "坐标化"为椭圆曲线上的一个点 $(x,y)$。

  \item \textbf{从曲线到环面}：给定 $\C$ 上任意一条椭圆曲线 $E$，
    存在格 $\Lambda$（在相似变换下唯一）使得
    $E(\C) \cong \C/\Lambda$。
\end{enumerate}
\end{theorem}

\begin{example}[正方形格对应的椭圆曲线]
\label{ex:square-lattice-curve}
取 $\Lambda = \Z + \Z i$（正方形格）。
由对称性 $\Lambda = i\Lambda$（乘以 $i$ 就是旋转 $90°$，
正方形格不变），可以证明 $g_3 = 0$。
对应的椭圆曲线是 $y^2 = 4x^3 - g_2 x$（没有常数项）。
这条曲线有一个额外的自同构 $(x,y) \mapsto (-x, iy)$，
对应于环面上的 $z \mapsto iz$。
\end{example}

\begin{proof}[Theorem~\ref{thm:uniformization} 的证明思路]
方向 (1) 的关键是利用 $\wp$ 的微分方程
$(\wp')^2 = 4\wp^3 - g_2\wp - g_3$，
说明 $(z \bmod \Lambda) \mapsto (\wp(z), \wp'(z))$
确实落在椭圆曲线上，并且是双射。
详细证明见\cref{thm:weierstrass-properties}。

方向 (2) 用到 $j$-不变量 $j(\tau)$ 给出双射
$\SL_2(\Z) \backslash \HH \to \C$，
所以给定任何椭圆曲线，总能找到对应的 $\tau$（从而找到格）。
\end{proof}

\textbf{同源：椭圆曲线之间的"好"映射。}
椭圆曲线之间最重要的映射叫做\textbf{同源}（isogeny）。
在费马大定理的证明中，Ribet 的关键步骤就涉及到一条椭圆曲线的"$3$-同源"——
他证明了如果 Frey 曲线存在，就能构造一个 $3$-同源导致矛盾。

\textbf{为什么同源比一般的映射更重要？}
因为同源\emph{保持群结构}（$\varphi(P+Q) = \varphi(P)+\varphi(Q)$），
所以它是"椭圆曲线的范畴"中的态射。
两条椭圆曲线之间如果存在同源，它们的 $L$-函数和 Galois 表示就密切相关——
实际上，同源的椭圆曲线在所有大素数处有相同的迹，
这是谷山-志村猜想中"$L$-函数相等"的核心。

\textbf{在复环面语言中，同源就是倍乘。}
$\C/\Lambda_1 \to \C/\Lambda_2$ 的同源对应 $\alpha \in \C^*$ 使得 $\alpha\Lambda_1 \subset \Lambda_2$，
映射是 $z \mapsto \alpha z$。最简单的情形是整数倍乘 $z \mapsto nz$，
核是 $E[n] = \{P \in E \mid nP = O\} \cong (\Z/n\Z)^2$（$n^2$ 个点）。

\begin{definition}[同源]
\label{def:isogeny}
\index{同源 (isogeny)}
设 $E_1, E_2$ 为椭圆曲线。一个\textbf{同源}（isogeny）
$\varphi\colon E_1 \to E_2$ 是保持群结构的非常数代数映射
（即 $\varphi(P + Q) = \varphi(P) + \varphi(Q)$）。

在复环面语言中，$\C/\Lambda_1$ 到 $\C/\Lambda_2$ 的同源对应于
$\alpha \in \C^*$ 使得 $\alpha\Lambda_1 \subset \Lambda_2$
（注意不要求相等，只要求包含），映射为 $z \mapsto \alpha z$。
\end{definition}

\begin{example}[乘法同源]
\label{ex:multiplication-isogeny}
取 $\Lambda = \Z + \Z i$，$\alpha = 2$。
则 $2\Lambda = 2\Z + 2\Z i \subset \Lambda$，
所以 $z \mapsto 2z$ 是 $\C/\Lambda$ 到自身的同源。
它的\textbf{度}是 $[\Lambda : 2\Lambda] = 4$
（因为 $\Lambda/2\Lambda \cong (\Z/2\Z)^2$ 有 $4$ 个元素），
意味着每个像有 $4$ 个原像（即核 $\Lambda/2\Lambda$ 有 $4$ 个元素：
$0, \; 1/2, \; i/2, \; (1+i)/2$）。
\end{example}


% ============================================================================

\section{模曲线与模空间}
\label{sec:modular-curves-moduli}
% ============================================================================

\textbf{为什么需要模曲线？}

\cref{cor:torus-tau} 告诉我们：$Y(1) = \SL_2(\Z) \backslash \HH$
参数化所有椭圆曲线的同构类，通过 $j$-不变量与 $\C$ 同构。
这很好——但有时我们不只想分类椭圆曲线，
还想分类"带有额外数据"的椭圆曲线。

\textbf{典型的"额外数据"是什么？}
以 Wiles 证明 FLT 为例。核心步骤是：把"级别 $N$ 的模形式"与"有理椭圆曲线"对应起来。
这里"级别 $N$"就意味着椭圆曲线上记录了某种 $N$ 阶结构（具体是一个 $N$ 阶循环子群）。
要参数化"带 $N$ 阶子群的椭圆曲线"，$Y(1)$ 就不够用了——它把有无额外结构的曲线混为一谈。

\textbf{解决方案：用更小的对称群。}
如果我们只用 $\Gamma_0(N)$ 的等价关系来识别椭圆曲线
（即只把"$N$ 阶子群相同"的曲线视为等价），
就得到一个更细的分类空间 $Y_0(N) = \Gamma_0(N) \backslash \HH$。
不同的 $N$ 和不同的子群类型（$\Gamma_0, \Gamma_1, \Gamma$）
给出不同"粗细程度"的分类。

这些分类空间——\textbf{模曲线}——不只是集合，
它们还具有代数曲线的结构（将在\cref{ch:riemann-surfaces}中看到这是紧黎曼面），
可以定义在有理数甚至整数上，成为研究椭圆曲线算术的核心工具。

\textbf{谷山-志村猜想}（现在是定理）断言：每条有理椭圆曲线 $E$ 都是某个模曲线 $X_0(N)$ 的"商"，
即存在有理映射 $X_0(N) \to E$。这就是 FLT 证明的骨架。

\begin{definition}[模曲线]
\label{def:modular-curves}
对同余子群 $\Gamma$，定义\textbf{（开）模曲线}
\[
  Y(\Gamma) = \Gamma \backslash \HH.
\]
特别地，$Y_0(N) = \GammaO(N) \backslash \HH$，
$Y_1(N) = \GammaI(N) \backslash \HH$，
$Y(N) = \Gamma(N) \backslash \HH$。

通过添加尖点得到\textbf{（紧化）模曲线}
\[
  X(\Gamma) = \Gamma \backslash \HH^*.
\]
$X(\Gamma)$ 是紧黎曼面（将在\cref{ch:riemann-surfaces}中详细讨论）。
\end{definition}

\begin{theorem}[模诠释]
\label{thm:moduli-interpretation}
模曲线参数化带附加结构的椭圆曲线（在 $\C$ 上）：
\begin{enumerate}
  \item $Y(1) = \SL_2(\Z) \backslash \HH$ 参数化复环面的同构类
    （\cref{cor:torus-tau}）。通过 $j$-不变量，$Y(1) \cong \C$。
  \item $Y_0(N)$ 参数化 $\{(E, C)\}$ 的同构类，其中 $E$ 是椭圆曲线，
    $C \subset E$ 是 $N$ 阶循环子群（等价于核为 $\Z/N\Z$ 的同源
    $E \to E/C$）。
  \item $Y_1(N)$ 参数化 $\{(E, P)\}$ 的同构类，其中 $P \in E$ 是
    $N$ 阶点。
  \item $Y(N)$ 参数化 $\{(E, (P, Q))\}$ 的同构类，其中 $(P, Q)$ 是
    $E[N]$ 的一组有序基。
\end{enumerate}
\end{theorem}

\begin{proof}[解释]
\textbf{(2)}：设 $\tau \in \HH$，对应环面 $E_\tau = \C/\Lambda_\tau$
（$\Lambda_\tau = \Z\tau + \Z$）。取 $C = \langle 1/N \rangle \subset E_\tau$，
即 $C = \frac{1}{N}\Z/\Z$ 是 $N$ 阶循环子群。

何时 $(E_\tau, C_\tau) \cong (E_{\tau'}, C_{\tau'})$？
需要 $\alpha \in \C^*$ 使得 $\alpha\Lambda_\tau = \Lambda_{\tau'}$
且 $\alpha C_\tau = C_{\tau'}$。写出条件：
$\alpha\tau = a\tau' + b$，$\alpha = c\tau' + d$（$ad - bc = 1$），
且 $\alpha \cdot \frac{1}{N} \equiv \frac{j}{N} \pmod{\Lambda_{\tau'}}$
对某 $j$ 与 $N$ 互素。

这等价于 $\frac{c\tau' + d}{N} \in \frac{1}{N}\Z + \Lambda_{\tau'}$，
即 $c \equiv 0 \pmod{N}$。
故等价关系正是 $\GammaO(N)$ 的作用。

\textbf{(3) 和 (4)}：类似分析。
\end{proof}

\subsection*{Hauptmodul 与亏格 $0$ 的函数域}

CSS 中反复强调：模曲线除了是参数空间，也有自己的函数域。对亏格 $0$ 的曲线，这个函数域最简单——它总是由一个函数生成。

\begin{definition}[Hauptmodul]
\label{def:hauptmodul}
\index{Hauptmodul}
设 $X$ 是定义在 $\C$ 上的紧黎曼面，且 $\genus(X)=0$。若一个亚纯函数
\[
  t\colon X \longrightarrow \mathbb{P}^1(\C)
\]
给出双有理同构，等价地说 $\C(X)=\C(t)$，则称 $t$ 为 $X$ 的\textbf{Hauptmodul}（主模）。
\end{definition}

\begin{proposition}[亏格 $0$ 曲线的函数域]
\label{prop:genus0-hauptmodul}
若 $X$ 是亏格 $0$ 的紧黎曼面，则 $X\cong \mathbb{P}^1(\C)$，从而存在 Hauptmodul $t$ 使得
\[
  \C(X)=\C(t).
\]

若 $P\in X$ 是一点，则可取 $t$ 在 $P$ 处有一个单极点、在别处全纯；这样的 $t$ 已经生成整个函数域。
\end{proposition}

\begin{proof}[证明思路]
由黎曼面分类，亏格 $0$ 的紧黎曼面与 $\mathbb{P}^1(\C)$ 双全纯同构。把该同构复合到坐标函数 $z$ 上，就得到一个亚纯函数 $t$。而 $\mathbb{P}^1$ 的亚纯函数正是有理函数，所以 $\C(X)$ 同构于 $\C(z)=\C(t)$。
\end{proof}

对模曲线来说，这个定义尤其重要：当 $X(\Gamma)$ 的亏格为 $0$ 时，我们就能用一个单参数来描述整条曲线，所有模函数都是这个参数的有理函数。

\begin{example}[$j$ 与若干低级别 Hauptmodul]
\label{ex:hauptmodul-examples}
\leavevmode
\begin{itemize}
  \item 对 $X(1)=\SL_2(\Z)\backslash \HH^*$，经典 $j$-不变量在尖点 $\infty$ 处有简单极点，且没有别的极点，因此
    \[
      \C(X(1))=\C(j).
    \]
    也就是说，$j(\tau)$ 正是 $X(1)$ 的 Hauptmodul。
  \item 对经典模曲线 $X(2)$，Hauptmodul 是 Legendre 的 $\lambda(\tau)$。它把带一组有序 $2$-扭点基的椭圆曲线写成
    \[
      y^2=x(x-1)(x-\lambda).
    \]
  \item 对 $X_0(2)$，一个常用的 Hauptmodul 是
    \[
      t_2(\tau)=\left(\frac{\eta(\tau)}{\eta(2\tau)}\right)^{24}=q^{-1}+O(1),
    \]
    其中 $\eta(\tau)=q^{1/24}\prod_{n\ge 1}(1-q^n)$ 是 Dedekind $\eta$ 函数。它满足
    \[
      j(\tau)=\frac{(t_2+256)^3}{t_2^2},
      \qquad
      j(2\tau)=\frac{(t_2+16)^3}{t_2}.
    \]
    因而 $j(\tau)$ 与 $j(2\tau)$ 都是 $t_2$ 的有理函数，说明 $\C(X_0(2))=\C(t_2)$。
  \item 对 $X_0(3)$，类似地可取
    \[
      t_3(\tau)=\left(\frac{\eta(\tau)}{\eta(3\tau)}\right)^{12}=q^{-1}+O(1),
    \]
    并有
    \[
      j(\tau)=\frac{(t_3+27)(t_3+243)^3}{t_3^3}.
    \]
    所以当 $X_0(N)$ 亏格为 $0$ 时，确实也存在一个单参数 Hauptmodul 来生成其函数域。
\end{itemize}
\end{example}

\subsection*{模点、同源类与扭曲}

前面已经从复环面角度说明了 $\SL_2(\Z)\backslash \HH$ 的意义。把这个结论改写成模曲线语言，就得到最基本的模诠释。

\begin{corollary}[复椭圆曲线的同构类]
\label{cor:complex-ec-moduli}
椭圆曲线 $E/\C$ 的同构类与 $Y(1)=\SL_2(\Z)\backslash \HH$ 的点一一对应；加入尖点后，$X(1)$ 只是把退化的结点三次曲线也补进去。

特别地，每个 $E/\C$ 都同构于某个 $E_\tau=\C/(\Z\tau+\Z)$，而且
\[
  E_\tau \cong E_{\tau'}
  \quad\Longleftrightarrow\quad
  \tau' = \gamma(\tau) \text{ 对某个 } \gamma\in \SL_2(\Z).
\]
\end{corollary}

\begin{proof}
这只是把 \cref{cor:torus-tau,thm:uniformization} 与模曲线记号合在一起重述。
\end{proof}

\begin{definition}[有理同源类]
\label{def:q-isogeny-class}
设 $E,E'/\Q$ 为椭圆曲线。若存在定义在 $\Q$ 上的同源 $E\to E'$，则称它们属于同一个\textbf{有理同源类}（isogeny class）。
\end{definition}

在模性定理里，真正与一个归一化新形式对应的通常不是单独一条椭圆曲线，而是一个 $\Q$-同源类：同一同源类中的曲线有相同的导子与同一个 Hasse--Weil $L$-函数，只是最小模型、挠点结构或 Tamagawa 数可能不同。

\begin{example}[两个特殊的 $j$ 值]
\label{ex:j0-j1728}
\leavevmode
\begin{itemize}
  \item $j=1728$：过 $\C$ 同构后，所有这样的椭圆曲线都可写成
    \[
      y^2=x^3-x,
    \]
    它有额外自同构 $z\mapsto iz$，所以
    \[
      \End(E_{\overline{\Q}})\supset \Z[i].
    \]
    过 $\Q$ 时，可以写成 $E_A\colon y^2=x^3+Ax$（$A\in \Q^\times$）；并且 $E_A\cong_{\Q}E_{A'}$ 当且仅当 $A'/A\in (\Q^\times)^4$。
  \item $j=0$：过 $\C$ 同构后，所有这样的椭圆曲线都可写成
    \[
      y^2=x^3-1,
    \]
    它有额外自同构 $z\mapsto \omega z$（$\omega=e^{2\pi i/3}$），所以
    \[
      \End(E_{\overline{\Q}})\supset \Z[\omega].
    \]
    过 $\Q$ 时，可以写成 $E_B\colon y^2=x^3+B$（$B\in \Q^\times$）；并且 $E_B\cong_{\Q}E_{B'}$ 当且仅当 $B'/B\in (\Q^\times)^6$。
\end{itemize}

这两个例子说明：同一个复同构类在 $\Q$ 上往往会裂成许多不同的有理同构类，而这些差别正是由扭曲控制的。
\end{example}

\begin{definition}[二次扭曲]
\label{def:quadratic-twist}
设
\[
  E\colon y^2=x^3+Ax+B
  \qquad (A,B\in \Q,\,4A^3+27B^2\neq 0)
\]
且 $d\in \Q^\times$。$E$ 关于 $d$ 的\textbf{二次扭曲}定义为
\[
  E^{(d)}\colon y^2=x^3+d^2Ax+d^3B.
\]
在 $\overline{\Q}$ 上，取 $\sqrt d$ 后映射
\[
  (x,y)\longmapsto (dx,d^{3/2}y)
\]
给出 $E\cong E^{(d)}$；但若 $d\notin (\Q^\times)^2$，则它们通常并不同构于 $\Q$。
\end{definition}

\textbf{注。} 扭曲的统一语言来自 Galois 上同调。固定 $E/\Q$ 后，所有在 $\overline{\Q}$ 上与 $E$ 同构、但在 $\Q$ 上可能不同构的模型，由集合
\[
  H^1\bigl(\Gal(\overline{\Q}/\Q),\Aut(E_{\overline{\Q}})\bigr)
\]
分类。若 $j(E)\neq 0,1728$，则 $\Aut(E_{\overline{\Q}})=\{\pm 1\}$，于是二次扭曲正对应于
\[
  H^1\bigl(\Gal(\overline{\Q}/\Q),\{\pm 1\}\bigr)
  \cong \Q^\times/(\Q^\times)^2.
\]
这就是为什么“按平方类扭曲”会自然出现。


% ============================================================================

\section{习题}
\label{sec:ch1-exercises}
% ============================================================================

第一章搭起了整本书的舞台：模群 $\SL_2(\Z)$ 的几何、同余子群的定义、模形式的变换律、椭圆曲线的基本性质。
这些概念表面上彼此独立，但它们其实在回答同一个问题——
如何用群的对称性来组织上半平面和格的算术。
下面的习题帮你把这条主线真正走一遍：
从具体的矩阵计算出发，一步步看清为什么模群的结构、基本域的形状、
模形式的权、以及椭圆曲线的等价关系，都是同一个对称故事的不同侧面。

做完这九题之后，你应当能够：
在脑子里快速判断任何两个格点是否 $\SL_2(\Z)$-等价；
用变换律推导出"奇权模形式只能是零"；
计算小级别同余子群的指标和尖点数；
并且理解 $j$-不变量为什么是椭圆曲线同构类的完全不变量。

% -----------------------------------------------------------------------
\begin{exercise}
\label{ex:ch1-sl2z-generators}
设 $S = \begin{pmatrix}0 & -1 \\ 1 & 0\end{pmatrix}$，$T = \begin{pmatrix}1 & 1 \\ 0 & 1\end{pmatrix}$。
\begin{enumerate}
  \item[(1)] 直接验证 $S^2 = -I$ 以及 $(ST)^3 = -I$。
  \item[(2)] 设 $\bar{S}, \bar{T}$ 是 $S, T$ 在 $\PSL_2(\Z) = \SL_2(\Z)/\{\pm I\}$ 中的像。
    由 (1) 推导 $\bar{S}^2 = 1$ 以及 $(\bar{S}\bar{T})^3 = 1$。
  \item[(3)] $\PSL_2(\Z)$ 的表现定理告诉我们这两个关系已经是完整表现，即
    $\PSL_2(\Z) \cong \Z/2\Z * \Z/3\Z$（自由积）。
    试举例说明，阶为 $2$ 的元素 $\bar{S}$ 与阶为 $3$ 的元素 $\bar{S}\bar{T}$
    生成的自由积有多"大"：写出 $\PSL_2(\Z)$ 中 $5$ 个不同的元素。
\end{enumerate}
\end{exercise}

\begin{proof}[解答]
(1) 直接矩阵乘法：
\[
  S^2 = \begin{pmatrix}0&-1\\1&0\end{pmatrix}^2
      = \begin{pmatrix}-1&0\\0&-1\end{pmatrix} = -I.
\]
$ST = \begin{pmatrix}0&-1\\1&0\end{pmatrix}\begin{pmatrix}1&1\\0&1\end{pmatrix}
= \begin{pmatrix}0&-1\\1&1\end{pmatrix}$，于是
\[
  (ST)^2 = \begin{pmatrix}0&-1\\1&1\end{pmatrix}^2
         = \begin{pmatrix}-1&-1\\1&0\end{pmatrix},
  \quad
  (ST)^3 = \begin{pmatrix}-1&-1\\1&0\end{pmatrix}\begin{pmatrix}0&-1\\1&1\end{pmatrix}
         = \begin{pmatrix}-1&0\\0&-1\end{pmatrix} = -I.
\]

(2) 在 $\PSL_2(\Z)$ 中取模 $\{\pm I\}$：
$\bar{S}^2 = \overline{S^2} = \overline{-I} = \bar{I}$（单位元），
$(\bar{S}\bar{T})^3 = \overline{(ST)^3} = \overline{-I} = \bar{I}$。

(3) 五个不同元素（取不同的字长）：
\[
  \bar{I},\quad \bar{S},\quad \bar{T},\quad \bar{S}\bar{T},\quad \bar{T}^2.
\]
因为 $\Z/2*\Z/3$ 是无限群，事实上 $\PSL_2(\Z)$ 也是无限群——光是 $\bar{T}^n$（$n \in \Z$）
就给出无穷多个不同元素。
\end{proof}

% -----------------------------------------------------------------------
\begin{exercise}
\label{ex:ch1-fund-domain-action}
回忆基本域 $\mathcal{F} = \{\tau \in \HH : |\tau| \geq 1,\; |\mathrm{Re}(\tau)| \leq \tfrac{1}{2}\}$。
\begin{enumerate}
  \item[(1)] 证明 $\mathcal{F}$ 中没有两个内点是 $\SL_2(\Z)$-等价的。
    （提示：若 $\gamma\tau = \tau'$，比较 $\mathrm{Im}(\tau)$ 和 $\mathrm{Im}(\tau')$；
    利用 $\mathrm{Im}(\gamma\tau) = \mathrm{Im}(\tau)/|c\tau+d|^2 \leq \mathrm{Im}(\tau)$
    对 $|c\tau+d|\geq 1$ 成立。）
  \item[(2)] 确定 $\mathcal{F}$ 的边界哪些点被等同：$\tau$ 与 $\tau+1$ 等同（通过 $T$），
    以及 $\tau$ 与 $-1/\tau$ 等同（通过 $S$）的点各在哪里？
\end{enumerate}
\end{exercise}

\begin{proof}[解答]
(1) 设 $\tau, \tau' \in \mathcal{F}$ 是两个内点，$\gamma = \begin{pmatrix}a&b\\c&d\end{pmatrix}$，
满足 $\gamma\tau = \tau'$。
由变换公式 $\mathrm{Im}(\gamma\tau) = \mathrm{Im}(\tau)/|c\tau+d|^2$，
若 $c \neq 0$，则 $|c\tau+d|^2 \geq (\mathrm{Im}(c\tau))^2 = c^2 \mathrm{Im}(\tau)^2$；
当 $\mathrm{Im}(\tau) \geq \sqrt{3}/2$（即 $\tau$ 在 $\mathcal{F}$ 内部时）且 $|c| \geq 1$，
有 $|c\tau+d| > 1$，所以 $\mathrm{Im}(\tau') < \mathrm{Im}(\tau)$。
对称地，$\mathrm{Im}(\tau) < \mathrm{Im}(\tau')$，矛盾。
若 $c = 0$，则 $\gamma = \pm\begin{pmatrix}1&n\\0&1\end{pmatrix}$，即 $\tau' = \tau+n$；
两个内点的实部都在 $(-1/2,1/2)$ 内，只有 $n = 0$ 才能相差整数，故 $\tau = \tau'$。

(2) 左边界 $\mathrm{Re}(\tau) = -1/2$ 与右边界 $\mathrm{Re}(\tau) = 1/2$ 上的点通过 $T$（$\tau \mapsto \tau+1$）等同。
下圆弧 $|\tau| = 1$，$\mathrm{Re}(\tau) \leq 0$ 上的点 $\tau$ 与 $-1/\tau \in$ 右半弧通过 $S$（$\tau \mapsto -1/\tau$）等同。
特殊点：$\tau = i$（满足 $S \cdot i = i$，固定点，阶 $2$），$\tau = e^{2\pi i/3}$（满足 $ST \cdot e^{2\pi i/3} = e^{2\pi i/3}$，固定点，阶 $3$）。
\end{proof}

% -----------------------------------------------------------------------
\begin{exercise}
\label{ex:ch1-weight-odd}
设 $f \in \Mk_k(\Gamma)$，其中 $-I \in \Gamma$（例如 $\Gamma = \SL_2(\Z)$）。
\begin{enumerate}
  \item[(1)] 利用 $-I \in \Gamma$ 以及模变换律 $f(\gamma\tau) = (c\tau+d)^k f(\tau)$，
    证明若 $k$ 为奇数，则 $f \equiv 0$。
  \item[(2)] 这说明对 $\GammaI(N)$（$N \geq 3$），奇权空间 $\Mk_k(\GammaI(N))$ 可以非零。
    为什么？（提示：$-I \notin \GammaI(N)$ 当 $N \geq 3$。）
\end{enumerate}
\end{exercise}

\begin{proof}[解答]
(1) 取 $\gamma = -I = \begin{pmatrix}-1&0\\0&-1\end{pmatrix} \in \Gamma$。
则 $(-I)\tau = \tau$（分子分母都乘 $-1$），且 $c\tau+d = -1$，故
$(c\tau+d)^k = (-1)^k$。
模变换律给出 $f(\tau) = f((-I)\tau) = (-1)^k f(\tau)$。
若 $k$ 为奇数，则 $(-1)^k = -1$，所以 $f(\tau) = -f(\tau)$，即 $f \equiv 0$。

(2) 当 $N \geq 3$ 时，$-I \notin \GammaI(N)$（因为 $-I \equiv -I_2 \not\equiv I_2 \pmod{N}$，
而 $\GammaI(N)$ 要求矩阵模 $N$ 余单位矩阵）。
因此论证在 (1) 中的第一步就失败了——$-I$ 不是允许的变换，我们无法推出 $f(\tau) = (-1)^k f(\tau)$。
于是对 $\GammaI(N)$，奇权模形式空间可以非零。
\end{proof}

% -----------------------------------------------------------------------
\begin{exercise}
\label{ex:ch1-index-gamma0}
\begin{enumerate}
  \item[(1)] 证明 $[\SL_2(\Z) : \GammaO(N)] = N \prod_{p \mid N}(1 + p^{-1})$。
    （提示：先证 $N = p^m$ 的情形，再用 CRT 拼合。）
  \item[(2)] 计算 $[\SL_2(\Z):\GammaO(11)]$ 和 $[\SL_2(\Z):\GammaO(12)]$。
  \item[(3)] 指标公式为何暗示了 $X_0(N)$ 的亏格与 $N$ 的大小有关？
\end{enumerate}
\end{exercise}

\begin{proof}[解答]
(1) 约化同态 $\SL_2(\Z) \to \SL_2(\Z/N\Z)$ 是满射（标准结论），
其核为 $\Gamma(N)$，故 $[\SL_2(\Z):\Gamma(N)] = |\SL_2(\Z/N\Z)|$。
而 $\GammaO(N)/\Gamma(N)$ 对应 $\SL_2(\Z/N\Z)$ 中下三角矩阵（$c \equiv 0 \pmod N$），
即形如 $\begin{pmatrix}a&b\\0&d\end{pmatrix}$ 且 $ad \equiv 1$，个数为 $N\prod_{p\mid N}(1-p^{-2}) \cdot \varphi(N)^{-1} \cdot \varphi(N) = \ldots$
（标准计算给出 $|\SL_2(\Z/N\Z)| / |\text{上三角子群}| = N\prod_{p\mid N}(1+p^{-1})$，见 Diamond-Shurman 第1章。）

(2) $[\SL_2(\Z):\GammaO(11)] = 11(1+1/11) = 12$。
$[\SL_2(\Z):\GammaO(12)] = 12(1+1/2)(1+1/3) = 12 \cdot 3/2 \cdot 4/3 = 24$。

(3) 亏格公式（后续章节）表明 $g(X_0(N))$ 大致为 $\mu/12 + O(\sqrt{N})$，
其中 $\mu = [\SL_2(\Z):\GammaO(N)]$。
$\mu \sim N \prod_{p\mid N}(1+p^{-1})$ 随 $N$ 增长，故亏格也随 $N$ 增长——
这正是 $N$ 越大，空间 $S_2(\GammaO(N))$ 越丰富的几何原因。
\end{proof}

% -----------------------------------------------------------------------
\begin{exercise}
\label{ex:ch1-cusps-count}
设 $N \geq 1$。证明 $\GammaO(N)$ 的尖点（即 $\mathbb{P}^1(\Q)$ 的等价类）个数为
\[
  \sum_{d \mid N} \varphi\!\left(\gcd\!\left(d,\,\tfrac{N}{d}\right)\right),
\]
其中 $\varphi$ 是 Euler 函数。

对 $N = 1, 2, 3, 4, 6, 11$ 分别计算尖点数，并与下方表格比对：

\begin{center}
\begin{tabular}{cccccccc}
\hline
$N$ & $1$ & $2$ & $3$ & $4$ & $6$ & $11$ \\
\hline
尖点数 & $1$ & $2$ & $2$ & $3$ & $4$ & $2$ \\
\hline
\end{tabular}
\end{center}
\end{exercise}

\begin{proof}[解答]
$\GammaO(N)$ 在 $\mathbb{P}^1(\Q)$ 上的轨道由以下参数化给出：
$r/s \in \mathbb{P}^1(\Q)$（$\gcd(r,s)=1$）属于哪个轨道，
由 $\gcd(s, N)$ 的值（即分母与 $N$ 的公因子）决定——不同的 $d = \gcd(s,N)$ 给出不同的轨道族。
精细分析（参见 D\&S 定理3.8.3）得出轨道数恰为题目中的公式。

逐一计算：
\begin{itemize}
  \item $N=1$：唯一因子 $d=1$，$\varphi(\gcd(1,1))=\varphi(1)=1$。尖点数 $= 1$。
  \item $N=2$：$d=1$：$\varphi(\gcd(1,2))=\varphi(1)=1$；$d=2$：$\varphi(\gcd(2,1))=\varphi(1)=1$。共 $2$。
  \item $N=3$：类似 $N=2$，得 $2$。
  \item $N=4$：$d \in \{1,2,4\}$：$\varphi(\gcd(1,4))=1$，$\varphi(\gcd(2,2))=\varphi(2)=1$，$\varphi(\gcd(4,1))=1$。共 $3$。
  \item $N=6$：$d \in \{1,2,3,6\}$：$1 + 1 + 1 + 1 = 4$。
  \item $N=11$（素数）：$d \in \{1,11\}$：$\varphi(1)+\varphi(1)=2$。共 $2$。
\end{itemize}
与表格完全一致。\qedhere
\end{proof}

% -----------------------------------------------------------------------
\begin{exercise}
\label{ex:ch1-surjection-slnz}
证明约化同态 $\pi_N \colon \SL_2(\Z) \to \SL_2(\Z/N\Z)$ 是满射。

（提示：
先对 $N = p$ 素数证明。
取 $\begin{pmatrix}a&b\\c&d\end{pmatrix} \in \SL_2(\Z/p\Z)$，分情形 $p \nmid c$ 和 $p \mid c$，构造原像。
一般情形用中国剩余定理：$N = p_1^{r_1} \cdots p_s^{r_s}$，
再用 $\SL_2(\Z/p^r\Z) \cong \SL_2(\Z/p\Z) \times \cdots$ 归约到素数情形。）
\end{exercise}

\begin{proof}[解答]
\textbf{素数情形 $N = p$}：
设 $A = \begin{pmatrix}a&b\\c&d\end{pmatrix} \in \SL_2(\Z/p\Z)$，$\det A = ad-bc \equiv 1 \pmod{p}$。

\textit{情形一：$p \nmid c$。}取整数 $c', d'$ 满足 $c' \equiv c, d' \equiv d \pmod{p}$，$\gcd(c',p) = 1$。
用 Bezout：存在 $a', b' \in \Z$ 使得 $a'd' - b'c' = 1$（因为 $\gcd(c',d') = 1$，由 $p \nmid c'$ 以及 $\gcd(\det A, p) = 1$）。
则 $\begin{pmatrix}a'&b'\\c'&d'\end{pmatrix} \in \SL_2(\Z)$ 映射到 $A$。

\textit{情形二：$p \mid c$，则 $p \nmid a$（否则 $\det A \equiv 0$）。}
$A$ 乘以某个初等矩阵（列变换）可化归到情形一。

\textbf{一般情形}：$N = \prod p_i^{r_i}$，由 CRT，$\Z/N\Z \cong \prod \Z/p_i^{r_i}\Z$，
故 $\SL_2(\Z/N\Z) \cong \prod \SL_2(\Z/p_i^{r_i}\Z)$。
对每个 $p^r$，提升定理（Hensel 引理的矩阵版本）给出 $\SL_2(\Z/p\Z) \cong \SL_2(\Z/p^r\Z)$（作为商）。
因此 $\pi_N$ 是满射。\qedhere
\end{proof}

% -----------------------------------------------------------------------
\begin{exercise}
\label{ex:ch1-j-invariant}
椭圆曲线 $E \colon y^2 = 4x^3 - g_2 x - g_3$（Weierstrass 形式）的 $j$-不变量定义为
\[
  j(E) = 1728 \cdot \frac{g_2^3}{g_2^3 - 27g_3^2}.
\]
\begin{enumerate}
  \item[(1)] 对格 $\Lambda_\tau = \Z\tau + \Z$（$\tau \in \HH$），
    我们有 $g_2(\Lambda_\tau) = 60 G_4(\tau)$ 和 $g_3(\Lambda_\tau) = 140 G_6(\tau)$，
    其中 $G_k$ 是 Eisenstein 级数。
    证明 $j(\tau) := j(\C/\Lambda_\tau)$ 满足模变换律 $j(\gamma\tau) = j(\tau)$ 对所有 $\gamma \in \SL_2(\Z)$。
  \item[(2)] 证明 $j(\tau) = q^{-1} + 744 + 196884q + \cdots$（$q = e^{2\pi i\tau}$），其中 $q^{-1}$ 项说明 $j$ 在尖点 $i\infty$ 处有一个极点。
    由此说明 $j \colon \mathcal{H}/\SL_2(\Z) \to \C$ 是双射（全纯同构 $X(1) \cong \mathbb{P}^1(\C)$）。
  \item[(3)] 验证 $j(i) = 1728$ 和 $j(e^{2\pi i/3}) = 0$。这两个特殊值有什么几何意义？
\end{enumerate}
\end{exercise}

\begin{proof}[解答]
(1) $g_2(\Lambda_\tau)$ 和 $g_3(\Lambda_\tau)$ 分别是 $\Lambda_\tau$ 上的同一函数，而格 $\Lambda_{\gamma\tau} = (c\tau+d)^{-1}\Lambda_\tau$（相差一个非零复数倍）。
格伸缩 $\Lambda \mapsto \lambda\Lambda$ 给出 $g_2(\lambda\Lambda) = \lambda^{-4}g_2(\Lambda)$，$g_3(\lambda\Lambda) = \lambda^{-6}g_3(\Lambda)$。
代入 $\lambda = (c\tau+d)^{-1}$：
\begin{align*}
  g_2(\Lambda_{\gamma\tau}) &= (c\tau+d)^4 g_2(\Lambda_\tau),\\
  g_3(\Lambda_{\gamma\tau}) &= (c\tau+d)^6 g_3(\Lambda_\tau).
\end{align*}
所以
\[
  j(\gamma\tau) = 1728\cdot\frac{g_2(\Lambda_{\gamma\tau})^3}{g_2(\Lambda_{\gamma\tau})^3 - 27g_3(\Lambda_{\gamma\tau})^2}
  = 1728\cdot\frac{(c\tau+d)^{12}g_2^3}{(c\tau+d)^{12}(g_2^3 - 27g_3^2)}
  = j(\tau).
\]

(2) 利用 $G_4(\tau) = \frac{2\zeta(4)}{(2\pi i)^4} + O(q)$ 及类似展开，
可以验证判别式 $\Delta = g_2^3 - 27g_3^2 = (2\pi)^{12}q\prod(1-q^n)^{24}$（稍后在 Eisenstein 章节详述），
所以 $\Delta$ 在 $q = 0$ 处有一阶零点，故 $j = 1728 g_2^3/\Delta$ 在 $q=0$ 处有一阶极点，
展开即得 $j = q^{-1} + 744 + 196884q + \cdots$。
$j$ 是 $\mathcal{H}/\SL_2(\Z) = X(1)$ 上的亚纯函数，在尖点有唯一极点（一阶），
由黎曼面上有理函数的次数等于极点个数，$j$ 的次数为 $1$，即 $j \colon X(1) \to \mathbb{P}^1(\C)$ 是双射。

(3) 在 $\tau = i$ 处：$S \cdot i = i$，即 $i$ 是 $S$ 的固定点。
$g_3(i) = 0$（因为晶格 $\Lambda_i = \Z i + \Z$ 有4重对称性 $i\Lambda_i = \Lambda_i$，而 $g_3$ 是权 $6$ 的量，$(i)^6 g_3 = g_3 \Rightarrow g_3 = 0$）。
故 $j(i) = 1728 \cdot g_2^3/g_2^3 = 1728$。

在 $\tau = \omega = e^{2\pi i/3}$ 处：$\omega$ 有3重对称性 $\omega\Lambda_\omega = \Lambda_\omega$，
类似地 $g_2(\omega) = 0$（权4量在 $(\omega)^4 = \omega$ 作用下得 $g_2 = \omega g_2$，故 $g_2=0$），
$j(\omega) = 0$。

几何意义：$j(i) = 1728$ 是具有额外自同构 $z \mapsto iz$（阶4，等价类中只有 $i$）的椭圆曲线；
$j(\omega) = 0$ 是具有自同构 $z \mapsto \omega z$（阶6，等价类中只有 $\omega$）的椭圆曲线。
\end{proof}

% -----------------------------------------------------------------------
\begin{exercise}
\label{ex:ch1-modular-form-graded-ring}
设 $E_4, E_6 \in \Mk_*(\SL_2(\Z))$ 是标准化 Eisenstein 级数（前几章将详细构造，
此处只用它们的基本性质：$E_4$ 是权 $4$ 模形式，$E_6$ 是权 $6$ 模形式，$\Delta = (E_4^3 - E_6^2)/1728$ 是权 $12$ 的尖点形式）。
\begin{enumerate}
  \item[(1)] 利用维数公式 $\dim\Mk_k(\SL_2(\Z)) = \lfloor k/12\rfloor + \epsilon_k$
    （$\epsilon_k = 0$ 若 $k \equiv 2 \pmod{12}$，否则 $\epsilon_k = 1$），
    列出 $k = 4, 6, 8, 10, 12, 14$ 时 $\dim\Mk_k$ 的值。
  \item[(2)] 说明为什么 $E_4^2 \in \Mk_8$ 以及 $E_4 E_6 \in \Mk_{10}$ 是这些空间的基底。
  \item[(3)] 证明 $\Delta \in \Sk_{12}(\SL_2(\Z))$ 是非零的，
    并由 $\dim\Sk_{12} = 1$ 推断 $\Sk_{12}$ 由 $\Delta$ 生成。
\end{enumerate}
\end{exercise}

\begin{proof}[解答]
(1)
\[
\begin{array}{c|cccccc}
k & 4 & 6 & 8 & 10 & 12 & 14 \\
\hline
\dim\Mk_k & 1 & 1 & 1 & 1 & 2 & 1
\end{array}
\]
（详细推导见第5章维数公式。）

(2) $E_4^2$ 有权 $4+4=8$，全纯，在尖点有界（因为 $E_4$ 在尖点有界），
所以 $E_4^2 \in \Mk_8$。$\dim\Mk_8 = 1$ 且 $E_4^2 \neq 0$（它的常数项为 $1^2 = 1 \neq 0$），
故 $E_4^2$ 是 $\Mk_8$ 的基底。类似地，$E_4 E_6 \in \Mk_{10}$，$\dim\Mk_{10}=1$，常数项 $1 \neq 0$，故是 $\Mk_{10}$ 的基底。

(3) $\Delta$ 的 $q$-展开为 $\Delta = q - 24q^2 + 252q^3 - \cdots$（Ramanujan $\tau$ 函数），
在 $q = 0$ 处消失，故 $\Delta$ 是尖点形式。又 $\Delta \neq 0$（首项系数 $1 \neq 0$）。
$\dim\Sk_{12} = \dim\Mk_{12} - 1 = 2 - 1 = 1$（减去非尖点的 Eisenstein 部分一维），
故 $\Sk_{12} = \C \cdot \Delta$。\qedhere
\end{proof}

% -----------------------------------------------------------------------
\begin{exercise}
\label{ex:ch1-isogeny-and-endo}
设 $\Lambda = \Z\tau + \Z \subset \C$，$E = \C/\Lambda$。
\begin{enumerate}
  \item[(1)] 证明：$\C$ 到 $\C$ 的映射 $z \mapsto \alpha z$（$\alpha \in \C^\times$）
    诱导出椭圆曲线之间的同源 $E \to \C/(\alpha\Lambda)$，当且仅当 $\alpha\Lambda \supset \Lambda$
    （即 $\alpha^{-1}\Lambda \subset \Lambda$，或等价地 $\alpha^{-1} \in \mathrm{End}(E) \otimes \Q$）。
  \item[(2)] 证明 $\mathrm{End}(E) = \Z$，当 $\tau$ 不是任何虚二次域中的代数整数时（即 $E$ 没有复乘）。
    （提示：若 $\alpha\Lambda \subset \Lambda$，写出 $\alpha\tau = a\tau+b$，$\alpha \cdot 1 = c\tau+d$，
    联立并消去 $\tau$，说明 $\alpha$ 满足整系数二次方程；若 $\tau \notin \overline{\Q}$ 或 $\tau$ 不是虚二次整数，则 $\alpha \in \Z$。）
  \item[(3)] 计算 $\mathrm{End}(E)$，当 $\tau = i$（即 $\Lambda = \Z i + \Z$，$E$ 是 Gaussian 整数格）。
    验证 $\alpha = i$ 给出 $E$ 到自身的四阶自同构。
\end{enumerate}
\end{exercise}

\begin{proof}[解答]
(1) $z \mapsto \alpha z$ 在 $\C$ 上是线性映射，它诱导 $E = \C/\Lambda \to \C/(\alpha\Lambda)$ 当且仅当
将 $\Lambda$ 中的点映到 $\alpha\Lambda$ 中，即 $\alpha\Lambda \subset \alpha\Lambda$（自动成立）。
但我们要的是群同态 $\C/\Lambda \to \C/\Lambda$（同一个椭圆曲线），
即需要 $\alpha\Lambda \subset \Lambda$。这正是 $\alpha \in \mathrm{End}(E) \otimes_\Z \Q$ 中整数部分的条件。

(2) 若 $\alpha\Lambda \subset \Lambda$，则 $\alpha\tau = a\tau + b$ 和 $\alpha \cdot 1 = c\tau + d$，
其中 $a,b,c,d \in \Z$。第二个方程：$\alpha = c\tau + d$。代入第一个：
$(c\tau+d)\tau = a\tau + b$，即 $c\tau^2 + (d-a)\tau - b = 0$。
若 $c = 0$，则 $\alpha = d \in \Z$。若 $c \neq 0$，则 $\tau$ 满足整系数二次方程，
即 $\tau$ 是某虚二次域中的代数整数（CM情形）。在非 CM 情形下 $c = 0$，故 $\alpha \in \Z$，$\mathrm{End}(E) = \Z$。

(3) $\Lambda = \Z i + \Z$。令 $\alpha = i$：$i \cdot (i) = -1 \in \Lambda$，$i \cdot 1 = i \in \Lambda$，
故 $i\Lambda \subset \Lambda$。映射 $z \mapsto iz$ 是 $E$ 的自同构（因为它是双射同源）。
它的阶是 $4$（$i^4 = 1$，$i^2 = -1 \neq 1$，$i \neq 1$）。
因此 $\mathrm{End}(\C/(\Z i+\Z)) \supseteq \Z[i]$（Gaussian 整数）；
事实上 $\mathrm{End}(E) \cong \Z[i]$（虚二次数域 $\Q(i)$ 的整数环）。\qedhere
\end{proof}

